% !TeX root = ../深度学习与大语言模型技术原理.tex
% 由主文档通过 \input 载入；请编译主文档。
\setcounter{chapter}{25} % 保留原书章号。
\chapter{长上下文的表示、训练与评测}
\label{ch:longcontext}

\section{窗口长度与有效使用能力}

上下文窗口规定模型一次能够接受的序列范围，实际可处理长度还受显存与执行实现限制。支持输入更多词元，不等于能稳定利用其中任意位置的信息。远距离依赖、重复干扰、位置分布变化以及输出预算，都会影响有效上下文能力。

长序列同时增加激活、注意力运算和 KV 存储。FlashAttention 减少中间矩阵搬运与存储，CP 分担各设备负担，GQA 或 KV 量化减少缓存，稀疏注意力减少连接；这些方法解决不同成本，也具有不同函数等价性。要判断系统瓶颈，应分别核算预填充、解码和存储，而不是只调整一个最大长度配置。

\section{位置插值与频率缩放}

设训练最大位置约为 $N_0$，希望扩展到 $sN_0$，$s>1$。位置插值（positional interpolation）将位置 $m$ 映为 $m/s$，使 RoPE 相位 $(m/s)\omega_j$ 回到原有范围\cite{pi}。相应地，两个相邻位置的相位间距也缩小，局部位置分辨率改变；避免大相位外推并不免费。

另一类方法修改 RoPE 基数或按频段调整频率，使高频与低频成分采用不同插值程度。YaRN 提供了结合频率处理与注意力尺度调整的具体方案\cite{yarn}。这些方法应视为修改模型的位置几何并可能配合继续训练，不能只凭命题~\ref{prop:rope-relative} 保证长程泛化。

若缩放因子依当前序列长度动态变化，已缓存键是否需要重算必须明确。否则同一注意力中的查询和历史键可能使用不同频率约定，破坏预期相对位置结构。实现时应把位置配置纳入模型和缓存版本。

\section{长上下文训练与数据构造}

长上下文训练需要真正包含长距离依赖的数据。随机把许多无关短文拼成很长序列，虽然满足长度要求，却未必产生利用远处证据的学习信号。可采用完整文档、跨段问答、代码仓库、长对话和带证据位置标注的任务，但要避免生成规则过于简单，导致模型只学模板。

渐进增加序列长度有助于控制训练成本，仍需保留短任务评测以检测退化。训练采样应覆盖不同证据位置、依赖距离和干扰强度。长样本减少每个批量的独立文档数，梯度方差与数据混合也随之改变，不能只把原学习率与批量配置原样迁移。

\section{从针检索到跨段综合}

针检索（needle-in-a-haystack）把特定事实放在长文本某处，再查询该事实，便于控制位置和长度。它测试明确线索检索，却不能代表多跳推理、跨文档比较或全局总结。应同时设置多个相关证据、相互矛盾文本、版本变化和无答案任务，检查模型是否只抓住局部关键词。

评测需系统改变总长度、证据位置和干扰密度，并测量准确率、引用支持、TTFT、ITL 和峰值显存。与短上下文 RAG 基线比较时，应计入检索成本和索引准备，但不能把不同可见证据条件下的结果直接解释成模型记忆能力差异。

